\section*{A Dichotomy (author Section 6)}
We first make Theorem 14 more concrete by identifying some interesting parameter settings. One natural setting is where the number of queries is polynomial and the advantage is inverse polynomial.
\paragraph{Theorem 29.} Suppose $L$ is computable and there is a randomized polynomial-time non-adaptive reduction $F$ from $L$ to $\omega(\log(|q|))$-additively approximating Kolmogorov complexity such that $F$ is polynomially honest and has inverse polynomial advantage. Then $L\in\mathrm{AM}\cap\mathrm{coAM}$.
\begin{proof}We apply Theorem 14 with $k_n^M=\operatorname{poly}(n)$, $\epsilon(n)=1/\operatorname{poly}(n)$, and $\alpha(n)=O(\log(n))$. Hence we have $\beta_n(q)=O(\log(n))$. Since $F$ is polynomially honest, we have that $\beta_n(q)=O(\log(q))$ for every $q\in Q_x^M$ for $x$ of length $n$. Hence by the conclusion of Theorem 14, we have $L\in\mathrm{AM}\cap\mathrm{coAM}$.\end{proof}
Another natural setting is where the number of queries and the advantage are both constant. In this case, we can get a conclusion from a smaller upper bound on the approximation gap.
\paragraph{Theorem 30.} Suppose $L$ is computable and there is a computable function $\eta=\omega(1)$ such that there is a randomized polynomial-time non-adaptive reduction $F$ from $L$ to $(2\log(K(q))+\eta(|q|))$-additively approximating Kolmogorov complexity such that $F$ makes $O(1)$ queries, is weakly honest and has advantage $\Omega(1)$. Then $L\in\mathrm{AM}\cap\mathrm{coAM}$.
\begin{proof}Let $\lambda:\mathbb N\to\mathbb N$ be an unbounded function such that $F$ is $\lambda$-honest. When applying Theorem 14, we choose $k_n^M=O(1)$ and $\epsilon(n)=\Omega(1)$. We also choose $\alpha$ appropriately so that $\beta_n(q)\ge2\log(K(q))+\eta(|q|)$---this is possible since $F$ is weakly honest, by ensuring $\alpha(\lambda^{-1})=o(\eta)$. Once more, applying Theorem 14 gives us the desired condlusion.\end{proof}
Theorem 29 supplies the first part of our dichotomy result. We require the following result of Hirahara [28] for the second part of our dichotomy. This is Theorem 6.3 in [28], where it is stated without mentioning the honesty property of the reduction. The proof of Theorem 6.3 in [28] does yield a polynomially honest reduction.
\paragraph{Theorem 31 ([28]).} For any constant $c\in\mathbb N$, there exists an $\mathrm{EXP}^{\mathrm{NP}}$ computable function $G=\{G_n\}$, where $G_n$ maps $n-c\log(n)$ bits to $n$ bits, such that for any $L\in\mathrm{NEXP}$, there is a randomized polynomial-time polynomially honest non-adaptive reduction from $L$ to distinguishing $G$ from uniform.
\paragraph{Corollary 32.} For every constant $c\in\mathbb N$ and any $L\in\mathrm{NEXP}$, there is a randomized polynomial-time non-adaptive reduction $F$ from $L$ to $c\log(|q|)$-additively approximating Kolmogorov complexity such that $F$ is polynomially honest and has constant advantage.
\begin{proof}Let $c\in\mathbb N$, and let $K'$ be any oracle that $c\log(|q|)$-additively approximates Kolmogorov complexity. Wlog we assume $c\ge3$---note that establishing the conclusion for a constant $c$ also establishes it for any constant $c'\le c$. Applying Theorem 31, we have that there is a computable function $G=\{G_n\}$, where $G_n$ maps $n-3c\log(n)$ bits to $n$ bits such that for any $L\in\mathrm{NEXP}$, there is a randomized polynomial-time polynomially honest non-adaptive reduction from $L$ to $D$, where $D$ is any language distinguishing $G$ from the uniform distribution. We define $D(q)=1$ if $K'(q)<|q|-1.5c\log(|q|)$. Note that any output of the generator $G$ can be described with $\log(n)+n-3c\log(n)+O(1)$ bits, and hence has Kolmogorov complexity $<n-2.5c\log(n)$ for large enough $n$. This implies that for every $n$-bit output $y$ of the generator $G_n$, $K'(y)<n-1.5c\log(n)$. On the other hand, with probability at least $1-1/n^{c/2}$ over $y\sim U_n$, we have that $K(y)\ge n-0.5c\log(n)$, and hence $K'(y)\ge n-1.5c\log(n)$. Thus $K'$ does distinguish the output of $G_n$ from uniform. Since $D'$ can be computed with a single deterministic query to $K'$ of the same length as the input, we have that for any $L\in\mathrm{NEXP}$, there is a randomized polynomial-time polynomially honest non-adaptive reduction from $L$ to $K'$. Since this holds for any $K'$ that $c\log(|q|)$-additively approximates Kolmogorov complexity, the result follows.\end{proof}
Theorem 29 together with Corollary 32 establishes the dichotomy result in Theorem 3.
